

%\section{Title for Section describing implementation of Environment and Policy Network}
\section{Methodology}\label{methodology}
\begin{figure}
    \centering
    \resizebox{0.5\textwidth}{!}{\includegraphics{framework.pdf}}
    \caption{Proposed Framework; ``S'' stands for State which contains the input claim and the path that is built starting from the head entity, ``R'' stands for the Reward which is used to update the agent's policy function at the end of the episode, ``a'' stands for the action which the Agent sends to the environment} %\textcolor{red}{ @Gustav, Kindly please explain the below sentence that you have written in terms like S stands for what, R stands for what...etc}}
    %\textcolor{olive}{Policy guided Beam Search extracts explainable paths from Knowledge Graph. Paths are used to arrive at veracity verdict. Because the information from the paths is varied, a voting mechanism is used to arrive at a consensus - the most likely tail entity between the paths. The verdict is then used to reward the agent for its performance, thus the classification accuracy is used to optimize the path extracting agent.
    %S stands for State, which contains the input claim and the path that is built starting from the head entity. R stands for the Reward which is used to update the Agent's Policy function at the end of the episode.
    %a stands for the action which the Agent sends to the environment - based on the State and action, the environment sends the Agent a new updated state, with the action (relation) and an entity that the action leads to, appended to the state/path.}
    \label{fig:framework}
\vspace{-5pt}
\end{figure}

This section provides an overview of the proposed framework for explainable fact-checking (Figure \ref{fig:framework}). %which involves extracting evidential paths and arriving at a veracity verdict. 
We first explain the path extraction process (Section \ref{evi})  succeeded by the veracity verdict (Section \ref{ver}), and the environment in Section \ref{env}.

\subsection{Evidential path extraction}
\label{evi}
The path extraction process follows multi-hop reasoning practices in which a policy-based agent sequentially extends an inference path. In contrast to traditional multi-hop reasoning tasks where one part of the claim is hidden from the agent, the proposed framework presents the agent with the entire claim and, in addition, the inference path, as it is being formed, as observations. We boost the quantity and versatility of extracted evidential paths by combining the policy-based agent with the beam search algorithm. 
 
\noindent \textbf{Policy-based agent: } The policy-based agent receives a claim as input and, in return, extracts a path from the KG. The path extraction process follows the pattern of a Markov decision process (MDP). The agent considers an observation, comprising the input claim and current path, to decide upon an action with which to extend the current path, after which the path is updated, and the agent receives an updated observation. The process is repeated until the expiration of a predetermined number of time steps.
The process results in a path that begins with the claim's head entity and bears, %\textcolor{red}{@Gustav, is the word "bears" here correct?} \textcolor{olive}{Yes, meaning 'is semantically similar to'}, 
after sufficient training, a semantic similarity to the claim. And thus, the resulting path can be used to judge upon the veracity of the input claim. As semantically similar, the tail entities of the claim and the evidential paths should match for a true claim and differ for a false claim.
The policy function maps the observation to a probability distribution over the action space which comprises all of the unique relations within the KG. The distribution is then filtered to only include relations that connect the current tail entity of the path. This is done to only include relations that can be used to extend the path. A \textit{self-loop} relation is included within the action space in case the agent decides to remain at the current tail entity. %\textcolor{red}{@Gustav, is the next sentences same as Section \ref{implement} Policy Network, if it is same, can we remove it?} \textcolor{olive}{The sentence contains info not contained in the \ref{implement} section - policy net is trained to optimize path based classification. What follows after, starting with "The policy function is parametrized..." is the same as \ref{implement} and can be removed, I think. The reason why it's here, because they teach at BA to repeat everything in the thesis...}
The policy function is parametrized by a neural network to maximize the expected reward for the path-based classification of the claim.
%The policy function is parametrized by a fully-connected neural network (NN). The NN is composed of one hidden layer with Rectified Linear Unit (ReLU) \fixme{why ReLU?} \textcolor{olive}{I cannot remember a particular reason and can't find any records of experimentation. Probably out of habit, the same going for Adam opimizer} serving as the activation function and an output layer normalized by a softmax function. A stochastic gradient-based method is used for optimization (Adam) \fixme{why Adam?} introduced in \cite{Kingma2014}.

\noindent\textbf{Beam search: }
The beam search algorithm allows us to increase the number of paths which are extracted during an episode. 
\begin{comment}The beam search is guided by a heuristic function that takes into account the probability distribution calculated by the agent's policy function and also the cosine similarity between the candidate entities which the agent chosen actions lead to.
\end{comment}
It is customary to guide the relation choice in KG reasoning tasks, but to leave the entity choice to random chance. In this work, %\textcolor{red}{@Gustav, as said above beam search is guided by heuristic and in next sentences also heuristic,check once if this is correct?} \textcolor{olive}{Changed it a bit and commented out the first heuristic mention, because it contained some false/dated info (the heuristic has nothing to do with cosine similarity), how's it look?} 
we propose a heuristic function that takes into account the fitness of the entities which the relations lead to by leveraging the KG scoring function \cite{trouillon2016complex}. %\textcolor{red}{@Gustav, Clarification required, is this scoring function the below metric?}. 
By doing this, the search algorithm takes into account how fitting the resulting entity is as a correct tail entity for the evidential path by scoring it as the tail entity in the input triple. As the result of the scoring function is less important in the early steps of path extension, this function is the base of an exponential function, where the step number is the exponent. This function is added to the probability calculated for a given action by the policy function and the sum is used to compare viable path extension options to each other in the beam search process. The final equation stands (Equation \ref{eq1}): %\fixme{we need to define the equation and its parameter, @Gustav, kindly please write what each parameter stands for}
%$$\pi(a|s, \theta)$$ \textcolor{olive}{is the output of the policy function for the action, given the state s and policy function parameters theta. The significance of this part in the equation is - how confident is the Policy network that the chosen action leads to the correct tail entity.}
%$$Re(< r, e_{s}, \bar{e}_c >)$$ 
%\textcolor{olive}{is the scoring function of ComplEx embedding. r, es and ec are all vector representations and the arrows around them signify a dot product of the three vectors signifies their dot product. Re means the real part of the dot product. This part of the equation answers the question: how likely is it that the triple (claim head, claim relation, candidate entity) belong to the Knowledge graph as a true triple. In the beginning of the path building, we do not expect to find the correct tail entity, but after a few steps, we expect to have reached the correct tail entity. To mirror this, I have made this into an exponential function, with the step of the episode being the exponent. (step starts at zero, thus in the beginning, all of the candidate entities have an equal shot at being chosen for beams, and the weight of the decision is put on the policy functions output.}
%$$^{step}$$

\begin{equation}
\label{eq1}
\pi(a|s, \theta) + Re(< r, e_{s}, \bar{e}_c >)^{step}
\end{equation}
%$$\pi(a|s, \theta) + Re(< r, e_{s}, \bar{e}_c >)^{step}$$
where, $\pi(a|s, \theta)$ is the output probability of the policy function for the action a, given state s and policy function parameters $\theta$, and $Re(< r, e_{s}, \bar{e}_c >)$ %\fixme{we have not defined the parameters in $Re(< r, e_{s}, \bar{e}_c >)$ } 
is the scoring function and $step$ is the step of the MDP episode. Here, $r$, $e_{s}$, and $\bar{e}_c$ are vector representations of relation, head entity, and tail entity. The tuple signifies the dot product of the three vectors. $Re$ means the real part of the dot product.


\subsection{Veracity verdict}\label{ver}
The veracity verdict is arrived upon by grouping the extracted paths and counting all unique path tail entities. The straightforward approach would be to choose the most popular tail entity as the choice tail entity to be compared to the claim tail, but since the result can be a tie, a heuristic weighting is added to the votes. The beam search heuristic for the final step of each path is used to weight the vote of each entity. This practically excludes the possibility of a tie between the contending tail entities and adds leverage to the entities which the agent is more confident in and which received a higher score from the KG triple scoring function which is used to evaluate whether a given triple belongs to the KG or not and is used to fit vector embeddings to entities and relations.
The verdict is then evaluated against the ground truth and the result is used to update the agent which extracted the paths.

\subsection{Environment}\label{env}

The environment is modeled as a Markov decision process (MDP) which controls the interaction between the RL agent (the Agent) and the knowledge graph. A 4-tuple $\mathcal{<S, A, T, R>}$ represents the MDP. Each element is described below.

\begin{comment}
\subsubsection{Episode}

In the beginning of an episode the agent receives an input triple that represents a fact claim. The agent then traverses the knowledge graph beginning with the input triple's source entity. This traversal produces a path that leads to the input triple's target entity if the fact claim is true or to some other entity if the claim is false or the agent is unable to find a correct path. This path should also provide as an explanation for the veracity verdict. Each path produced has a span of 4 entities connected as a chain by 3 relations. A new relation-entity pair is appended to the path at every step.
A training episode ends with the agent being rewarded. A testing episode ends with a vote between found paths to produce a veracity verdict for the input triple.

\end{comment}

\noindent \textbf{State ($\mathcal{S}$):} The state consists of entities and relations. Two representations of the state are operated with in parallel: 
1) a human-readable collection of real-world elements: entities and relations, and 
2) a corresponding collection of embeddings - vector representations that capture the semantic information of the state elements. The state contains the fact claim, traversed path, and current location of the agent which is itself an entity. The initial state consists of an input triple which is the fact claim. %This triple is reversed so that, \textcolor{red}{@Gustav, do we need to give the below example or we can remove it?}
%$$\mathcal<Joe Biden, presidentOf, USA>$$
%$$\mathcal<USA, presidentOf, Joe Biden>.$$ 
Thus the initial state can be thought of constituting as, %\fixme{can we do something about the following state? we can make it part of previous line ..?}
%$$\mathcal<claimed Target, claimed Relation, beginningOfPath, ...>.$$
$<$ \textit{claimed Target}, \textit{claimed Relation}, \textit{beginningOfPath}, ...$>$

The embedded state is padded with zero vectors since the state is used as input to the policy network which requires a fixed-size input. These zero vectors are incrementally replaced by relation-entity embedding pairs during the agent's path traversal process. The final entity in the set of these pairs is of particular importance: this is the entity that is used to give a veracity verdict for the fact claim. So in essence, once the path traversal is completed then the veracity verdict stands as: $$claimedTarget == finalEntity$$

\noindent\textbf{Action ($\mathcal{A}$): }
Not only is the input size of the policy network fixed but so is the output size - it is always the same size as the action space. The policy network takes the state as input and computes probabilities for each action in the action space. This can be thought of as the agent making an estimate: which actions are most likely to take me from the current entity to the desired final entity?
During training time, one of the top three most probable actions are chosen for each transition. The choice is random with a distribution that is calculated from the output of the policy network to fit the new sample size. %\textcolor{red}{@Gustav, what is this 3, is this required or redundant, if its redundant kindly remove it}(3). 
Random choice is introduced in order to reduce the likelihood of the agent getting stuck looking for sub-optimal paths brought forth by the local minima problem. %In other words, the random choice is introduced to encourage the Agent to explore a bit more.

\noindent\textbf{Transition ($\mathcal{T}$): }
Choosing a new action entails a transition of state. The transition itself is a function $T:$ $S X A$ $->$ $S$, which means that performing an action at a given state produces a new state. This new state is the input triple, a traversed path plus the choice action and a new current entity. The new current entity is a node in the KG that is connected to the previous current entity by the choice action. It is not uncommon for an entity to have multiple neighbors that are all connected to it via the same relation. %For example, the entity \textit{The Beatles} will have 4 neighbors \fixme{why we expect reader to assume 4 neighbors? Should we put a footnote about Beatles had 4 group members ? or it is fine?} connected to it via the relation \textit{memberOf}.
\begin{comment}\textcolor{red}{@Gustav, is the next sentence required or will it be better to remove it?}
\end{comment}
Most previous literature concerning multi-hop reasoning picks a neighboring entity at random \cite{Xiong2017, Das2017}, the focus being on the performance of the RL agent, unsuited for the task of choosing a candidate from a variable size set of entities.
\begin{comment}
\textcolor{red}{@Gustav, is this same as Beam Search heuristic above, if it is same, can we remove the next sentence?}In this work a heuristic is proposed for choosing a most likely entity to extend the path in the following subsection "Beam search and voting".
\end{comment}

\noindent\textbf{Reward ($\mathcal{R}$): }
The agent is given a reward once the path traversal is complete. The reward is 1 if the final entity in the found path is equal to the true target.
Otherwise, the reward is 0.
\[
    reward = 
\begin{cases}
    1,& \text{if final entity = true target}\\
    0,              & \text{otherwise}
\end{cases}
\]

\begin{comment}
\textcolor{red}{@Gustav, do we really need the below example or can we remove it, I feel like it is redundant as the above condition is sufficient?}
For example, if the fact claim is $$<Joe Biden, presidentOf, Russia>$$ then the Agent is given a reward of 1 if it arrives at \textit{USA} and a reward 0 for anything else.
\end{comment}